% Part: incompleteness
% Chapter: theories-computability
% Section: intepretability

\documentclass[../../../include/open-logic-section]{subfiles}

\begin{document}

\olfileid{inc}{tcp}{itp}

\olsection{$\Th{Q}$ను అర్థనిర్దేశం చేయగల సిద్ధాంతాలు అనిర్ణయనీయాలు}

ఈ ఫలితాలను ఇంకా బలపరచవచ్చు. అనౌపచారికంగా చెప్పాలంటే, ఒక భాష
$\Lang{L_1}$ను మరొక భాష~$\Lang{L_2}$లో అర్థనిర్దేశం చేయడంలో, $\Lang{L_1}$
యొక్క విశ్వాన్ని, సంబంధ సంకేతాలను, ప్రమేయ సంకేతాలను $\Lang{L_2}$లోని
\tetoken{సూత్రాలతో}{!!{formula}s} నిర్వచించడం ఉంటుంది. ఇక్కడ దానికి సమయం కేటాయించకపోయినా, ఈ
నిర్వచనాన్ని ఖచ్చితంగా చెప్పవచ్చు.

\begin{thm}
  అంకగణిత భాషను అర్థనిర్దేశం చేయగల భాషలోని సిద్ధాంతం~$\Th{T}$ను
  పరిగణించండి; ఆ అర్థనిర్దేశంలోని $\Th{Q}$తో $\Th{T}$ అవైరుధ్యంగా ఉందని
  అనుకుందాం. అప్పుడు $\Th{T}$ అనిర్ణయనీయం. $\Th{Q}$ స్వీకృతాల
  అర్థనిర్దేశాన్ని $\Th{T}$ నిరూపిస్తే, $\Th{T}$ యొక్క ఏ అవైరుధ్య విస్తరణా
  నిర్ణయించదగినది కాదు.
\end{thm}

ఈ నిరూపణ ముందటి సిద్ధాంత నిరూపణకు చేసిన చిన్న మార్పు మాత్రమే; ఒక
ప్రతివాద ఉదాహరణను ఉపయోగించి $\Th{Q}$, $\Th{\bar Q}$ల వేరుపరిచే సమితిని
పొందవచ్చు. ఎంపిక స్వీకృతంతో కూడిన జెర్మెలో--ఫ్రెంకెల్ సమితి సిద్ధాంతం
$\Th{ZFC}$ను, సాధారణ గణితంలో చాలావరకు నిర్వహించగలిగేంత శక్తిమంతమైన
స్వీకృతాధార పునాదిగా తీసుకోవచ్చు. $\Th{ZFC}$లో సహజ సంఖ్యలను నిర్వచించవచ్చు;
ఈ అర్థనిర్దేశం ద్వారా $\Th{Q}$ స్వీకృతాలు సత్యమవుతాయి. కాబట్టి మనకు

\begin{cor}
$\Th{ZFC}$కు అవైరుధ్యమైన, నిర్ణయించదగిన విస్తరణ ఏదీ లేదు.\sourcecorrection{OLTETCP-003}{ఆంగ్ల మూలం ``decidable extension'' అని మాత్రమే చెప్పి అవైరుధ్య షరతును వదిలేసింది. కానీ ముందరి సిద్ధాంతం అవైరుధ్య విస్తరణలకే వర్తిస్తుంది; అంతేకాక వైరుధ్య విస్తరణలో ప్రతి వాక్యమూ నిరూపించదగినది కనుక అది నిర్ణయించదగిన ప్రతివాద ఉదాహరణ. అందువల్ల అవసరమైన ``అవైరుధ్యమైన'' షరతును పునరుద్ధరించాం.}
\end{cor}

\begin{cor}
$\Th{ZFC}$కు సంపూర్ణమైన, అవైరుధ్యమైన, గణనీయంగా
\tetoken{స్వీకృతీకరించదగిన}{!!{axiomatizable}} విస్తరణ ఏదీ లేదు.
\end{cor}

$\Th{ZFC}$ భాషలో ఒకే ద్విస్థానిక సంబంధం~$\in$ మాత్రమే ఉంది. (నిజానికి
సర్వసమానత కూడా అవసరం లేదు.) కాబట్టి మనకు

\begin{cor}
ద్విస్థానిక సంబంధ సంకేతం గల ఏ భాషకైనా మొదటిస్థాయి తర్కం అనిర్ణయనీయం.
\end{cor}

ఈ ఫలితం రెండు ఏకస్థానిక ప్రమేయ సంకేతాలు గల ప్రతి భాషకూ విస్తరిస్తుంది;
ఎందుకంటే వాటితో ఒక ద్విస్థానిక సంబంధ సంకేతాన్ని అనుకరించవచ్చు. ఇప్పుడే
చెప్పిన ఫలితాలు నిశితమైనవి: కేవలం \emph{ఏకస్థానిక} సంబంధ సంకేతాలు, అత్యధికంగా
ఒక \emph{ఏకస్థానిక} ప్రమేయ సంకేతం మాత్రమే గల భాషకు మొదటిస్థాయి తర్కం
నిర్ణయించదగినదని తేలుతుంది.

ఇంకొక చిన్న విశేషం. ప్రామాణిక నమూనాలో సత్యమైన $\Obj 0$, $'$, $+$,
$\times$, $<$ భాషలోని వాక్యాల సమితి అనిర్ణయనీయమని మనకు తెలుసు. నిజానికి
మిగిలిన సంకేతాల ద్వారా $<$ను నిర్వచించవచ్చు; ఆ తరువాత $\times$, $'$ల ద్వారా
$+$ను నిర్వచించవచ్చు. కాబట్టి $\Obj 0$, $'$, $\times$ భాషలోని సత్య
వాక్యాల సమితి అనిర్ణయనీయం. మరోవైపు, $\Obj 0$, $'$, $+$ భాషలో అంకగణితంలో
సత్యమైన వాక్యాల సమితి నిర్ణయించదగినదని ప్రెస్‌బర్గర్ చూపాడు. అయితే ఆ పద్ధతి
గణనాత్మకంగా ఆచరణసాధ్యం కాదు.

\end{document}
